\section*{Введение}
\addcontentsline{toc}{section}{Введение}

В отчёте представлено описание выполнения курсовой работы по дисциплине <<Теория графов>>. В работе требуется разработать экспертную систему, основанную на бинарном дереве решений и продукционной модели представления знаний.

Разрабатываемая система предназначена для выбора параметров шрифтовой гарнитуры при подготовке бейджей к офсетной печати. Пользователь отвечает на вопросы о составе макета, алфавите, условиях чтения и технологических особенностях печати, после чего система выдаёт набор рекомендаций по выбору гарнитуры, насыщенности, контраста и межбуквенного интервала.

Разработанная консольная программа на языке C++ объединяет результаты пяти лабораторных работ. Программа выполняет случайную генерацию графа и весовых матриц, вычисляет эксцентриситеты вершин, выделяет центр графа и диаметральные вершины, ищет маршруты, кратчайшие пути, точки сочленения и потоки, строит минимальный остов, кодирует его кодом Прюфера, находит максимальное независимое множество рёбер, строит эйлеров цикл и фундаментальную систему разрезов. Исходный граф, сформированный в первой работе, затем используется в последующих работах.

В отчёте рассмотрены математическое описание использованных алгоритмов, особенности их программной реализации и результаты работы программы на сгенерированных графах.
